\documentclass[aspectratio=169, 12pt]{beamer}
\usepackage{graphicx}
\usepackage{hyperref}
\hypersetup{
    colorlinks=true,
    urlcolor=MidnightBlue,
    linkcolor=MidnightBlue,
    citecolor=MidnightBlue
}
\usepackage{booktabs}
\usepackage{enumerate}
\usepackage[absolute,overlay]{textpos}
\usepackage{amsfonts}
\usepackage{colortbl}
\usepackage{tcolorbox}
\usepackage[default]{lato}
\usepackage[dvipsnames]{xcolor}
\usepackage{pgfplots}
\pgfplotsset{compat=1.16}

\setbeamertemplate{itemize items}[default]
\setbeamertemplate{section in toc}[circle]
\setbeamertemplate{subsection in toc}[circle]
\setbeamertemplate{navigation symbols}{}
\setbeamertemplate{caption}[numbered]
\usepackage{comment}
\usepackage{textcomp}

\newcommand{\backupbegin}{
   \newcounter{finalframe}
   \setcounter{finalframe}{\value{framenumber}}
}
\newcommand{\backupend}{
   \setcounter{framenumber}{\value{finalframe}}
}

\AtBeginSection[]{
  \begin{frame}
  \vfill
  \centering
    \usebeamerfont{title}\usebeamercolor[fg]{title} \LARGE \insertsectionhead\par%
  \vfill
  \end{frame}
}

\newcommand{\hlcode}[1]{\tcbox[colback=orange!10, colframe=orange!10, arc=3pt, boxrule=0pt, left=2pt, right=2pt, top=1pt, bottom=1pt, on line]{\texttt{#1}}}

\setbeamerfont{frametitle}{size=\large}
\setbeamertemplate{footline}[frame number]
\setbeamersize{text margin left=4mm, text margin right=4mm}

\begin{document}
\linespread{1.3}

\title{Session 2: AI for Coding --- Practical Workflows and Demos}
\author[AI for Economists]{Elira Kuka and Tanner Regan}
\date[\today]{The George Washington University}
\maketitle


%%%%%%%%%%%%%%
\begin{frame}
\frametitle{Plan for Today: hands-on workflows}

\begin{enumerate}
    \item Coding with Claude Code overview
    \begin{itemize}
        \item Writing and debugging code from plain English {\color{gray}\it EK}
        \item ``Vibe coding'': building a complete tool without writing a line of code {\color{gray}\it EK}
    \end{itemize}

    \item Data analysis demos with Stata
    \begin{itemize}
        \item Stata example with Claude Code from terminal {\color{gray}\it EK}
        \item Stata example with Claude Code from VS Code {\color{gray}\it TR}
    \end{itemize}

    \item Python in VS Code overview
    \begin{itemize}
        \item libraries and environments {\color{gray}\it TR}
        \item Jupyter extension {\color{gray}\it TR}
    \end{itemize}

    \item Data analysis demo with Python
    \begin{itemize}
        \item Python example with Claude Code from VS Code {\color{gray}\it TR}
    \end{itemize}
\end{enumerate}

\end{frame}


\section{\textcolor{violet}{Coding with Claude Code overview \vspace{0.3in}}}


%%%%%%%%%%%%%%
\begin{frame}
\frametitle{Writing Code from Plain English}

\begin{itemize}
    \item You describe the task; AI writes the code --- in any language (Stata, R, Python, Julia)
    \item This works even if you do not know the language
    \item Key skill: writing a good prompt
        \begin{itemize}
            \item[$\circ$] Describe inputs and outputs concretely: ``I have a panel dataset with vars X, Y, Z\ldots''
            \item[$\circ$] Specify what you want verified: ``run it and show me the first 5 rows of output''
            \item[$\circ$] Be explicit about conventions: ``use `i.year` for year fixed effects in Stata''
        \end{itemize}
    \item \textcolor{gray}{\textit{$\triangleright$ Placeholder: side-by-side example of prompt $\to$ code}}
\end{itemize}

\end{frame}


%%%%%%%%%%%%%%
\begin{frame}
\frametitle{Debugging Code with AI}

\begin{itemize}
    \item Paste the error message and ask: ``what is wrong and how do I fix it?''
    \item With Claude Code in agentic mode: it sees the error itself and iterates automatically
    \item Particularly powerful for:
        \begin{itemize}
            \item[$\circ$] Obscure Stata/R errors with cryptic messages
            \item[$\circ$] Merge problems (wrong key, many-to-many, unmatched observations)
            \item[$\circ$] Output that looks right but is subtly wrong (off-by-one, wrong unit of observation)
        \end{itemize}
    \vspace{.05in}
    \item Important caveat: AI can fix the error while introducing a different bug
        \begin{itemize}
            \item[$\circ$] Always verify output against known benchmarks or summary statistics
        \end{itemize}
\end{itemize}

\end{frame}


%%%%%%%%%%%%%%
\begin{frame}
\frametitle{``Vibe Coding'': Building Tools Without Writing Code}

\begin{itemize}
    \item Concept: describe what you want at a high level; AI writes, runs, and refines all the code
    \item You act as the product manager, not the programmer
    \item Examples relevant to economists:
        \begin{itemize}
            \item[$\circ$] ``Build a script that downloads monthly CPI from FRED, merges with my wage data, and plots real wages by year''
            \item[$\circ$] ``Create a do-file that runs my DiD specification with all the standard robustness checks and exports tables to LaTeX''
            \item[$\circ$] ``Write a web scraper that pulls contract award data from USASpending.gov''
        \end{itemize}
    \item When it works well; when to be cautious
\end{itemize}

\end{frame}


%%%%%%%%%%%%%%
\begin{frame}
\frametitle{Automating Repetitive Data Tasks}

\begin{itemize}
    \item Common tasks AI handles easily:
        \begin{itemize}
            \item[$\circ$] Cleaning and reshaping messy data (wide $\leftrightarrow$ long, renaming, recoding)
            \item[$\circ$] Labeling variables and values from a codebook
            \item[$\circ$] Merging multiple datasets on complex keys
            \item[$\circ$] Generating a full set of summary statistics or balance tables
        \end{itemize}
    \vspace{.05in}
    \item Practical workflow for cleaning:
        \begin{enumerate}
            \item Show Claude the raw data structure (variable list, codebook, first rows)
            \item Describe the cleaning steps in plain English
            \item Review the output, ask for changes, iterate
        \end{enumerate}
    \item \textcolor{gray}{\textit{$\triangleright$ Placeholder: example or screenshot}}
\end{itemize}

\end{frame}


%%%%%%%%%%%%%%
\begin{frame}
\frametitle{Recap: Three Modes of Operation}

\begin{itemize}
    \item \textbf{Ask before edits} (safest):
        \begin{itemize}
            \item[$\circ$] Claude proposes each file edit before making it; you approve or reject
            \item[$\circ$] Best for: getting started, when you want full control, unfamiliar codebases
        \end{itemize}
    \vspace{.05in}
    \item \textbf{Plan mode} (recommended for most tasks):
        \begin{itemize}
            \item[$\circ$] Claude writes a full plan in Markdown; you review and approve before anything runs
            \item[$\circ$] Produces clean, commented code with an explicit audit trail
            \item[$\circ$] Can view and revisit session history
        \end{itemize}
    \vspace{.05in}
    \item \textbf{Auto mode} (most powerful, most risk):
        \begin{itemize}
            \item[$\circ$] Claude acts autonomously until the task is done
            \item[$\circ$] Use for well-defined tasks; always check the output
        \end{itemize}
\end{itemize}

\end{frame}


%%%%%%%%%%%%%%
\begin{frame}
\frametitle{Demo: Writing a Data Cleaning Script (Ask Mode)}

\begin{itemize}
    \item \textcolor{gray}{\textit{$\triangleright$ Demo slide}}
    \item Task: take a raw dataset, produce a clean analysis file
    \item Steps we will walk through:
        \begin{enumerate}
            \item Point Claude Code at the project folder
            \item Write a minimum working example in ``Ask before edits'' mode
            \item Claude proposes each change; we approve
            \item Verify output matches expectations
        \end{enumerate}
    \item What to watch for:
        \begin{itemize}
            \item[$\circ$] Does the code run without errors?
            \item[$\circ$] Do summary statistics match what we expect?
            \item[$\circ$] Any implicit assumptions we need to check?
        \end{itemize}
\end{itemize}

\end{frame}


%%%%%%%%%%%%%%
\begin{frame}
\frametitle{Demo: Cleaning Script Continued (Plan Mode)}

\begin{itemize}
    \item \textcolor{gray}{\textit{$\triangleright$ Demo slide}}
    \item Now run the same task in ``Plan mode'':
        \begin{enumerate}
            \item Ask Claude to write a plan before touching any files
            \item Review the plan in Markdown; request changes to the approach
            \item Approve; Claude executes and produces a clean, commented script
            \item Compare output to Ask-mode version: same result, better code
        \end{enumerate}
    \item Reviewing session history: how to audit what Claude did
    \item Ask Claude to update README with what was done (good practice)
\end{itemize}

\end{frame}


%%%%%%%%%%%%%%
\begin{frame}
\frametitle{When to Trust AI Code --- and How to Verify It}

\begin{itemize}
    \item AI code can look completely correct and be wrong in subtle ways
    \item Verification checklist:
        \begin{itemize}
            \item[$\circ$] Does the code run without errors? (necessary, not sufficient)
            \item[$\circ$] Do observation counts make sense at each merge/collapse step?
            \item[$\circ$] Do means and distributions match codebook or prior literature?
            \item[$\circ$] Can you reproduce a known result (e.g., a published Table 1)?
            \item[$\circ$] Ask Claude: ``what assumptions did you make?'' --- it will often reveal them
        \end{itemize}
    \vspace{.05in}
    \item General rule: trust AI code for \textit{structure}; verify AI code for \textit{content}
\end{itemize}

\end{frame}


%%%%%%%%%%%%%%
\begin{frame}
\frametitle{Building a Replication Package with Claude Code}

\begin{itemize}
    \item Claude Code can read your entire project and produce a replication package:
        \begin{enumerate}
            \item Point Claude at the project directory
            \item Ask it to trace the data pipeline from raw files to final tables
            \item Ask it to write a master script that runs everything in order
            \item Ask it to write a README documenting data sources, software, and steps
        \end{enumerate}
    \item This task would normally take days --- Claude does it in minutes
    \item You still need to verify: does running the master script reproduce the tables?
    \item AEA replication standards: what the package needs to include
\end{itemize}

\end{frame}


\section{\textcolor{violet}{Data analysis demos with Stata \vspace{0.3in}}}


%%%%%%%%%%%%%%
\begin{frame}
\frametitle{Using Stata with Claude Code}

\begin{itemize}
    \item \textcolor{gray}{\textit{$\triangleright$ Demo slide --- Elira}}
    \item Claude Code works with Stata via the terminal or VS Code (Stata Enhanced extension)
    \item What you can ask Claude to do:
        \begin{itemize}
            \item[$\circ$] Write a do-file from a plain-English description of the analysis
            \item[$\circ$] Debug a Stata error (paste the error; Claude explains and fixes it)
            \item[$\circ$] Translate a complex cleaning routine into clean, commented Stata code
            \item[$\circ$] Export regression tables to LaTeX (\texttt{esttab}, \texttt{outreg2})
        \end{itemize}
    \vspace{.05in}
    \item Example: [PLACEHOLDER --- data analysis example to demonstrate]
\end{itemize}

\end{frame}


%%%%%%%%%%%%%%
\begin{frame}
\frametitle{Stata Demo: Data Analysis with Claude Code}

\begin{itemize}
    \item \textcolor{gray}{\textit{$\triangleright$ Demo slide}}
    \item Task: \textcolor{gray}{\textit{[Placeholder: insert specific empirical example]}}
    \item Walk-through:
        \begin{enumerate}
            \item Describe the dataset and goal to Claude in plain English
            \item Claude writes the do-file; we run it
            \item Inspect output, ask for revisions
            \item Ask Claude to add comments and export a summary table
        \end{enumerate}
\end{itemize}

\end{frame}


\section{\textcolor{violet}{Python in VS Code overview \vspace{0.3in}}}

%%%%%%%%%%%%%%
\begin{frame}
\frametitle{Why Python?}

\begin{itemize}
    \item 
\end{itemize}
Python for Economists: Common Tasks
\begin{itemize}
    \item \textcolor{gray}{\textit{$\triangleright$ Placeholder: example or screenshot}}
    \item Data analysis with \texttt{pandas}: merging, collapsing, reshaping --- analogous to Stata
    \item Econometrics: \texttt{statsmodels}, \texttt{linearmodels} (panel data, IV)
    \item Data collection: \texttt{requests}, \texttt{BeautifulSoup} for scraping (Session 3)
    \item Figures: \texttt{matplotlib}, \texttt{seaborn} --- publication-quality plots
    \vspace{.05in}
    \item Claude Code can translate your Stata code to Python (or vice versa) --- useful for:
        \begin{itemize}
            \item[$\circ$] Running analyses you cannot do in Stata
            \item[$\circ$] Working with large datasets where Python is faster
        \end{itemize}
\end{itemize}

\end{frame}


%%%%%%%%%%%%%%
\begin{frame}
\frametitle{Python basics}

\begin{itemize}
    \item 
\end{itemize}

\end{frame}

%%%%%%%%%%%%%%
\begin{frame}
\frametitle{Libraries and Environments}

\begin{itemize}
    \item \textcolor{gray}{\textit{$\triangleright$ Demo slide --- Tanner}}
    \item Anaconda manages Python installations and packages
    \item Each project gets its own \textbf{environment}: a fixed set of packages and versions
        \begin{itemize}
            \item[$\circ$] Prevents conflicts between projects
            \item[$\circ$] Ensures others can reproduce your results exactly
        \end{itemize}
    \item Ask Claude Code to create and activate an environment for you
\end{itemize}
%https://claesbackman.com/claude-code-guide.html#:~:text=of%20Economics.-,Python,-ms%2Dpython.python
\end{frame}


%%%%%%%%%%%%%%
\begin{frame}
\frametitle{Python Setup in VS Code}

\begin{itemize}
    \item \textcolor{gray}{\textit{$\triangleright$ Demo slide --- Tanner}}
    \item Required extensions: Python, Jupyter
    \item Managing environments:
        \begin{itemize}
            \item[$\circ$] Each project gets its own environment (set of packages + versions)
            \item[$\circ$] Prevents conflicts between projects
            \item[$\circ$] Ask Claude Code to create and activate an environment for you
        \end{itemize}
    \item Running code:
        \begin{itemize}
            \item[$\circ$] \textbf{Jupyter notebook} (\texttt{.ipynb}): interactive cells, inline output --- good for exploration
            \item[$\circ$] \textbf{Jupyter interactive window}: run \texttt{.py} scripts cell-by-cell --- best of both worlds
            \item[$\circ$] \textbf{Script} (\texttt{.py}): run the full pipeline --- best for production code
        \end{itemize}
\end{itemize}

\end{frame}


%%%%%%%%%%%%%%
\begin{frame}
\frametitle{Running Python Code with Jupyter}

\begin{itemize}
    \item Three modes in VS Code:
        \begin{itemize}
            \item[$\circ$] \textbf{Jupyter notebook} (\texttt{.ipynb}): interactive cells, inline output --- good for exploration
            \item[$\circ$] \textbf{Interactive window}: run \texttt{.py} scripts cell-by-cell --- best of both worlds
            \item[$\circ$] \textbf{Script} (\texttt{.py}): run the full pipeline --- best for production code
        \end{itemize}
    \item Claude Code can write, run, and debug Python in any of these modes
\end{itemize}

\end{frame}


\section{Data Analysis Demo with Python \vspace{0.3in}}

%%%%%%%%%%%%%%
\begin{frame}
\frametitle{Python Exercise}

\begin{itemize}
    \item Download data from Kigali City
    \item Download data from Alpha Earth
    \item Visualize source data
    \item Build a sample frame
    \item Visualize the sample frame
\end{itemize}

\end{frame}

%%%%%%%%%%%%%%
\begin{frame}
\frametitle{Download data from Kigali City}

\begin{enumerate}
    \item Wrote out the prompt "Prompt data scraping for Kigali.md"
    \item Claude generated the plan \href{https://github.com/tannerregan/GWU_AI_study_group/blob/main/code/lecture_2/python_demo/kigali-city-data-mutable-frost.md}{"kigali-city-data-mutable-frost.md"} ($\sim$10 mins)
    \item Ran it without edits (using Jupyter window). Two issues: 
    \begin{itemize}
        \item Claude added packages to the yml but didn't update the conda environment. Chatting with Claude we identified this andfixed it. Told claude to update the project CLAUDE.md file so that it always: (1) updates the python environment with newly required packages, and (2) updates the .yml file so that the environment is perfectly replicable.
        \item I had to leave for daycare pickup and the code was still running. 
- I told Claude "the script is running. I need to close my laptop and go home. Can I pause the script and restart it without disrupting the download?"
-CC: "The script has no resume logic — if you stop and restart it, it will re-download everything from scratch, including layers already completed. The fix is simple: add a file-exists check so it skips layers that are already on disk. Want me to add that now before you leave? It's a two-line change per download function and would make the restart seamless."
- CC updated the script so I could go home. SHOW: GitHub Diff between the two files.
    \end{itemize}
    \item Ran the updated script at home and the rest downloaded
\end{enumerate}

\end{frame}




\section{\textcolor{violet}{Wrapping Up \vspace{0.3in}}}


%%%%%%%%%%%%%%
\begin{frame}
\frametitle{Key Takeaways from Sessions 1 \& 2}

\begin{enumerate}
    \item AI handles the coding mechanics --- your job is to verify the output and drive the analysis
    \smallskip
    \item Use Plan mode for any non-trivial task; it gives you an audit trail and catches mistakes early
    \smallskip
    \item Git + Claude Code = version-controlled, reproducible research without learning Git commands
    \smallskip
    \item CLAUDE.md files carry your preferences across sessions so you do not repeat yourself
    \smallskip
    \item Replication packages are now a hours-long task, not a weeks-long one
    \smallskip
    \item The critical skill is \textit{verification}: check outputs, not just that the code runs
\end{enumerate}

\end{frame}


%%%%%%%%%%%%%%
\begin{frame}
\frametitle{Things to Try Before the Next Session}

\begin{itemize}
    \item Set up VS Code with Claude Code, Git, and Python (use the setup guide)
    \item Ask Claude to write a small data task in your preferred language --- verify the output
    \item Initialize a Git repo for one of your current projects
    \item Write a CLAUDE.md file for that project
    \vspace{.05in}
    \item Coming up in Session 3: Finding and Collecting Data
        \begin{itemize}
            \item[$\circ$] Identifying datasets, scraping, APIs, parsing PDFs at scale
        \end{itemize}
\end{itemize}

\end{frame}



%%%%%%%%%%%%%%
\begin{frame}
\frametitle{Questions?}

\begin{itemize}
    \item Questions on anything from Sessions 1 or 2?
    \item What tasks in your own research would you most like to try this on?
\end{itemize}

\end{frame}


\end{document}
